\begin{thebibliography}{99}
\addcontentsline{toc}{section}{\refname}
\small
\bibitem{Aczel}
P.~Aczel, \emph{An introduction to inductive definitions},
in J.~Barwise (ed.), Handbook of Mathematical Logic,
North-Holland, 1977, pp.~739--782.
\href{https://lawrencecpaulson.github.io/papers/Aczel-Inductive-Defs.pdf}{Text}.

\bibitem{BauerHanson}
A.~Bauer and J.~E. Hanson, \emph{The Countable Reals},
\href{https://arxiv.org/abs/2404.01256v3}{arXiv:2404.01256v3}, 2026.

\bibitem{Bernardi}
C.~Bernardi, \emph{Precomplete equivalence relations in dominical
categories}, Proc. Amer. Math. Soc. \textbf{98} (1986), 124--134.

\bibitem{CockettLists}
J.~R.~B. Cockett, \emph{List-arithmetic distributive categories: Locoi},
J. Pure Appl. Algebra \textbf{66} (1990), no.~1, 1--29.
\href{https://doi.org/10.1016/0022-4049(90)90121-W}{doi:10.1016/0022-4049(90)90121-W}.

\bibitem{CockettComputability}
J.~R.~B. Cockett, \emph{Categories and Computability},
lecture notes, 12 August 2014.
\href{https://cspages.ucalgary.ca/~robin/talks/estonia-winter-2010/estonia-notes.pdf}{Author's text}.

\bibitem{CGHRangeI}
J.~R.~B. Cockett, X.~Guo, and P.~Hofstra,
\emph{Range categories I: General theory},
Theory and Applications of Categories \textbf{26} (2012),
no.~17, pp.~412--452.
\href{https://www.tac.mta.ca/tac/volumes/26/17/26-17abs.html}{Journal page}.

\bibitem{CGHRangeII}
J.~R.~B. Cockett, X.~Guo, and P.~Hofstra,
\emph{Range categories II: Towards regularity},
Theory and Applications of Categories \textbf{26} (2012),
no.~18, pp.~453--500.
\href{https://www.tac.mta.ca/tac/volumes/26/18/26-18abs.html}{Journal page}.

\bibitem{CockettHofstra}
J.~R.~B. Cockett and P.~J.~W. Hofstra,
\emph{Introduction to Turing categories},
Ann. Pure Appl. Logic \textbf{156} (2008), 183--209.

\bibitem{CockettLack}
J.~R.~B. Cockett and S.~Lack, \emph{Restriction categories. I.
Categories of partial maps}, Theoret. Comput. Sci. \textbf{270}
(2002), 223--259.

\bibitem{CockettLackIII}
J.~R.~B. Cockett and S.~Lack,
\emph{Restriction categories III: colimits, partial limits and extensivity},
Math. Structures Comput. Sci. \textbf{17} (2007), 775--817.

\bibitem{CurryCosheaves}
J.~M. Curry, \emph{Dualities between cellular sheaves and cosheaves},
J. Pure Appl. Algebra \textbf{222} (2018), 966--993.
\href{https://doi.org/10.1016/j.jpaa.2017.06.001}
{doi:10.1016/j.jpaa.2017.06.001}.

\bibitem{DPH}
R.~A. Di Paola and A.~Heller, \emph{Dominical categories:
recursion theory without elements}, J. Symbolic Logic
\textbf{52} (1987), no.~3, 594--635.
\href{https://doi.org/10.2307/2274352}{doi:10.2307/2274352}.

\bibitem{DPM}
R.~A. Di Paola and F.~Montagna, \emph{Some properties of the syntactic
$p$-recursion categories generated by consistent, recursively enumerable
extensions of Peano arithmetic}, J. Symbolic Logic \textbf{56} (1991),
643--660.

\bibitem{FPT}
M.~Fiore, G.~Plotkin, and D.~Turi,
\emph{Abstract syntax and variable binding},
Proceedings of the Fourteenth Annual IEEE Symposium on Logic in
Computer Science, IEEE Computer Society, 1999, pp.~193--202.
\href{https://doi.org/10.1109/LICS.1999.782615}{doi:10.1109/LICS.1999.782615}.

\bibitem{GabbayPitts}
M.~J. Gabbay and A.~M. Pitts,
\emph{A new approach to abstract syntax with variable binding},
Formal Aspects of Computing \textbf{13} (2002), 341--363.
\href{https://doi.org/10.1007/s001650200016}{doi:10.1007/s001650200016}.

\bibitem{GoguenBurstall}
J.~A. Goguen and R.~M. Burstall,
\emph{A study in the foundations of programming methodology:
Specifications, institutions, charters and parchments},
in D.~Pitt, S.~Abramsky, A.~Poign\'e, and D.~Rydeheard (eds.),
Category Theory and Computer Programming, Lecture Notes in Computer
Science \textbf{240}, Springer, 1986, pp.~313--333.
\href{https://doi.org/10.1007/3-540-17162-2_131}{doi:10.1007/3-540-17162-2\_131}.

\bibitem{GTWW}
J.~A. Goguen, J.~W. Thatcher, E.~G. Wagner, and J.~B. Wright,
\emph{Initial algebra semantics and continuous algebras},
J. ACM \textbf{24} (1977), no.~1, 68--95.
\href{https://doi.org/10.1145/321992.321997}{doi:10.1145/321992.321997}.

\bibitem{HP}
P.~H\'ajek and P.~Pudl\'ak, \emph{Metamathematics of First-Order
Arithmetic}, Perspectives in Logic, Cambridge \mbox{University} Press,
2017 reissue of the Springer text, originally published in 1993.
\href{https://doi.org/10.1017/9781316717271}{doi:10.1017/9781316717271}.

\bibitem{Heller}
A.~Heller, \emph{An existence theorem for recursion categories},
J. Symbolic Logic \textbf{55} (1990), 1252--1268.

\bibitem{HH}
H.~P.~Hirst and J.~L.~Hirst, \emph{A Primer for Logic and Proof},
2015 edition.
\href{https://www.appstate.edu/~hirstjl/primer/hirst.pdf}{Authors' text}.

\bibitem{Katsumata}
S.-y. Katsumata,
\emph{A generalisation of pre-logical predicates to simply typed formal systems},
in J.~D\'iaz, J.~Karhum\"aki, A.~Lepist\"o, and D.~Sannella (eds.),
Automata, Languages and Programming, Lecture Notes in Computer Science
\textbf{3142}, Springer, 2004, pp.~831--845.
\href{https://doi.org/10.1007/978-3-540-27836-8_70}{doi:10.1007/978-3-540-27836-8\_70}.

\bibitem{Keisler}
H.~J. Keisler, \emph{The ultraproduct construction},
author's survey, 1 June 2010.
\href{https://people.math.wisc.edu/~hkeisler/ultraproducts-web-final.pdf}{Author's text}.

\bibitem{KirstPeters}
D.~Kirst and B.~Peters, \emph{G\"odel's theorem without tears ---
essential incompleteness in synthetic computability},
in B.~Klin and E.~Pimentel (eds.), 31st EACSL Annual Conference on
Computer Science Logic (CSL 2023), Leibniz International Proceedings
in Informatics \textbf{252}, Schloss Dagstuhl, 2023, pp.~30:1--30:18.
\href{https://doi.org/10.4230/LIPIcs.CSL.2023.30}{doi:10.4230/LIPIcs.CSL.2023.30}.

\bibitem{Lengyel2004}
F.~Lengyel, \emph{More existence theorems for recursion categories},
Ann. Pure Appl. Logic \textbf{125} (2004), 1--41.

\bibitem{Lengyel2006}
F.~Lengyel, \emph{Locally connected recursion categories},
Technical Report TR-2006011, Graduate Center, City University of New York,
2006. \url{https://academicworks.cuny.edu/gc_cs_tr/277/}.

\bibitem{Corrections}
F.~Lengyel, \emph{Corrections to More existence theorems for recursion
categories}, version 10, 20 October 2026,
\href{https://github.com/flengyel/errata/blob/5555e4fa1da7b8e3a0836dd1fbf1ca385cacd03b/lengyel2004/lengyel2004_repairs_v10.tex}{source}.

\bibitem{Composition}
F.~Lengyel, \emph{Pointwise provable equality and the failure of
composition}, 2026,
\href{https://arxiv.org/abs/2609.25556}{arXiv:2609.25556}.

\bibitem{MaiettiLists}
M.~E. Maietti, \emph{Joyal's arithmetic universe as list-arithmetic pretopos},
Theory and Applications of Categories \textbf{24} (2010),
no.~3, 39--83.
\href{https://www.math.unipd.it/~maietti/papers/jartacrev.pdf}{Author's text}.

\bibitem{Miller2004}
J.~S. Miller, \emph{Degrees of unsolvability of continuous functions},
J. Symbolic Logic \textbf{69} (2004), no.~2, 555--584.
\href{https://doi.org/10.2178/jsl/1082418543}{doi:10.2178/jsl/1082418543}.

\bibitem{Montagna}
F.~Montagna, \emph{``Pathologies'' in two syntactic categories of partial
maps}, Notre Dame J. Formal Logic \textbf{30} (1989), 105--116.

\bibitem{MPV}
J.~Murwanashyaka, F.~Pakhomov, and A.~Visser,
\emph{There are no minimal essentially undecidable theories},
J. Logic Comput. \textbf{34} (2024), 1159--1171.
\href{https://doi.org/10.1093/logcom/exad005}{doi:10.1093/logcom/exad005}.

\bibitem{PourEl}
M.~B. Pour-El, \emph{Effectively extensible theories},
J. Symbolic Logic \textbf{33} (1968), no.~1, 56--68.
\href{https://doi.org/10.2307/2270052}{doi:10.2307/2270052}.

\bibitem{PourEl1970}
M.~B. Pour-El, \emph{A recursion-theoretic view of axiomatizable theories},
Dialectica \textbf{24} (1970), no.~4, 267--276.
\href{https://www.jstor.org/stable/42970364}{JSTOR 42970364}.
\href{https://pi.math.cornell.edu/~yytran/PourEl_recursion_axiomatizable_theories.pdf}{Full text}.

\bibitem{RobinsonRosolini}
E.~Robinson and G.~Rosolini, \emph{Categories of partial maps},
Inform. and Comput. \textbf{79} (1988), no.~2, 95--130.
\href{https://doi.org/10.1016/0890-5401(88)90034-X}{doi:10.1016/0890-5401(88)90034-X}.

\bibitem{RosoliniThesis}
G.~Rosolini, \emph{Continuity and effectiveness in topoi}, D.Phil.
thesis, University of Oxford, 1986.

\bibitem{Rosolini}
G.~Rosolini, \emph{Representation theorems for $P$-categories},
in F.~Borceux (ed.), \emph{Categorical Algebra and its Applications},
Lecture Notes in Mathematics \textbf{1348}, Springer, Berlin, 1988,
307--315.

\bibitem{vanDijkGO}
J.~van Dijk and A.~Gietelink Oldenziel,
\emph{G\"odel incompleteness through arithmetic universes after A.~Joyal},
2020. \href{https://arxiv.org/abs/2004.10482}{arXiv:2004.10482}.

\bibitem{VisserEHU}
A.~Visser, \emph{Essential hereditary undecidability},
Arch. Math. Logic \textbf{63} (2024), nos.~5--6, 529--562.
\href{https://doi.org/10.1007/s00153-024-00911-y}{doi:10.1007/s00153-024-00911-y}.

\bibitem{Voutsadakis}
G.~Voutsadakis, \emph{Categorical abstract algebraic logic:
categorical algebraization of first-order logic without terms},
Arch. Math. Logic \textbf{44} (2005), no.~4, 473--491.
\href{https://doi.org/10.1007/s00153-004-0266-7}{doi:10.1007/s00153-004-0266-7}.

\bibitem{Wissmann}
T.~Wi\ss{}mann, \emph{Minimality notions via factorization systems and examples},
Log. Methods Comput. Sci. \textbf{18} (2022),
no.~3, article~31, 31:1--31:26.
\href{https://doi.org/10.46298/lmcs-18(3:31)2022}{doi:10.46298/lmcs-18(3:31)2022}.
\end{thebibliography}
